\documentclass[10pt,a4paper]{amsart}
\usepackage{amsmath,amssymb,amsthm,graphicx}
\usepackage{verbatim}
\newenvironment{reference}{\comment}{\endcomment}
%\newenvironment{reference}{}{}
\newenvironment{slogan}{\comment}{\endcomment}
\newenvironment{history}{\comment}{\endcomment}

% For commutative diagrams we use Xy-pic
\usepackage[all]{xy}

% We use 2cell for 2-commutative diagrams.
\xyoption{2cell}
% Theorem environments.
%
\theoremstyle{plain}
\newtheorem{theorem}[subsection]{Theorem}
\newtheorem{proposition}[subsection]{Proposition}
\newtheorem{lemma}[subsection]{Lemma}

\theoremstyle{definition}
\newtheorem{definition}[subsection]{Definition}
\newtheorem{example}[subsection]{Example}
\newtheorem{exercise}[subsection]{Exercise}
\newtheorem{situation}[subsection]{Situation}

\theoremstyle{remark}
\newtheorem{remark}[subsection]{Remark}
\newtheorem{remarks}[subsection]{Remarks}

\numberwithin{equation}{subsection}

% Macros
%
\def\lim{\mathop{\mathrm{lim}}\nolimits}
\def\colim{\mathop{\mathrm{colim}}\nolimits}
\def\Spec{\mathop{\mathrm{Spec}}}
\def\Hom{\mathop{\mathrm{Hom}}\nolimits}
\def\Ext{\mathop{\mathrm{Ext}}\nolimits}
\def\SheafHom{\mathop{\mathcal{H}\!\mathit{om}}\nolimits}
\def\SheafExt{\mathop{\mathcal{E}\!\mathit{xt}}\nolimits}
\def\Sch{\mathit{Sch}}
\def\Mor{\mathop{\mathrm{Mor}}\nolimits}
\def\Ob{\mathop{\mathrm{Ob}}\nolimits}
\def\Sh{\mathop{\mathit{Sh}}\nolimits}
\def\NL{\mathop{N\!L}\nolimits}
\def\CH{\mathop{\mathrm{CH}}\nolimits}
\def\proetale{{pro\text{-}\acute{e}tale}}
\def\etale{{\acute{e}tale}}
\def\QCoh{\mathit{QCoh}}
\def\Ker{\mathop{\mathrm{Ker}}}
\def\Im{\mathop{\mathrm{Im}}}
\def\Coker{\mathop{\mathrm{Coker}}}
\def\Coim{\mathop{\mathrm{Coim}}}

% Boxtimes
%
\DeclareMathSymbol{\boxtimes}{\mathbin}{AMSa}{"02}

%
% Macros for moduli stacks/spaces
%
\def\QCohstack{\mathcal{QC}\!\mathit{oh}}
\def\Cohstack{\mathcal{C}\!\mathit{oh}}
\def\Spacesstack{\mathcal{S}\!\mathit{paces}}
\def\Quotfunctor{\mathrm{Quot}}
\def\Hilbfunctor{\mathrm{Hilb}}
\def\Curvesstack{\mathcal{C}\!\mathit{urves}}
\def\Polarizedstack{\mathcal{P}\!\mathit{olarized}}
\def\Complexesstack{\mathcal{C}\!\mathit{omplexes}}
% \Pic is the operator that assigns to X its picard group, usage \Pic(X)
% \Picardstack_{X/B} denotes the Picard stack of X over B
% \Picardfunctor_{X/B} denotes the Picard functor of X over B
\def\Pic{\mathop{\mathrm{Pic}}\nolimits}
\def\Picardstack{\mathcal{P}\!\mathit{ic}}
\def\Picardfunctor{\mathrm{Pic}}
\def\Deformationcategory{\mathcal{D}\!\mathit{ef}}
\usepackage{cleveref}
\crefname{originaldefinition}{original definition}{original definitions}
\Crefname{originaldefinition}{Original definition}{Original definitions}
\makeatletter
\newlabel{definition-point}{{00Y5}{source}{}{sourcepointer-1}{}}
\newlabel{definition-point@cref}{{[originaldefinition][0][]00Y5}{[1][0][]source}{}{}{}}
\makeatother
\crefname{originaldefinition}{original definition}{original definitions}
\Crefname{originaldefinition}{Original definition}{Original definitions}
\makeatletter
\newlabel{definition-topos}{{00XA}{source}{}{sourcepointer-2}{}}
\newlabel{definition-topos@cref}{{[originaldefinition][0][]00XA}{[1][0][]source}{}{}{}}
\makeatother
\crefname{originallemma}{original lemma}{original lemmas}
\Crefname{originallemma}{Original lemma}{Original lemmas}
\makeatletter
\newlabel{lemma-covering-surjective-after-sheafification}{{00WT}{source}{}{sourcepointer-3}{}}
\newlabel{lemma-covering-surjective-after-sheafification@cref}{{[originallemma][0][]00WT}{[1][0][]source}{}{}{}}
\makeatother
\crefname{originallemma}{original lemma}{original lemmas}
\Crefname{originallemma}{Original lemma}{Original lemmas}
\makeatletter
\newlabel{lemma-sheafification-exact}{{00WJ}{source}{}{sourcepointer-4}{}}
\newlabel{lemma-sheafification-exact@cref}{{[originallemma][0][]00WJ}{[1][0][]source}{}{}{}}
\makeatother
\crefname{originalequation}{original equation}{original equations}
\Crefname{originalequation}{Original equation}{Original equations}
\makeatletter
\newlabel{equation-stalk}{{04EH}{source}{}{sourcepointer-5}{}}
\newlabel{equation-stalk@cref}{{[originalequation][0][]04EH}{[1][0][]source}{}{}{}}
\makeatother
\crefname{originallemma}{original lemma}{original lemmas}
\Crefname{originallemma}{Original lemma}{Original lemmas}
\makeatletter
\newlabel{lemma-points-recover}{{00Y6}{source}{}{sourcepointer-6}{}}
\newlabel{lemma-points-recover@cref}{{[originallemma][0][]00Y6}{[1][0][]source}{}{}{}}
\makeatother
\crefname{originallemma}{original lemma}{original lemmas}
\Crefname{originallemma}{Original lemma}{Original lemmas}
\makeatletter
\newlabel{lemma-point-pushforward-sheaf}{{00Y8}{source}{}{sourcepointer-7}{}}
\newlabel{lemma-point-pushforward-sheaf@cref}{{[originallemma][0][]00Y8}{[1][0][]source}{}{}{}}
\makeatother
\crefname{originallemma}{original lemma}{original lemmas}
\Crefname{originallemma}{Original lemma}{Original lemmas}
\makeatletter
\newlabel{lemma-sheaf-coequalizer-representable}{{00WS}{source}{}{sourcepointer-8}{}}
\newlabel{lemma-sheaf-coequalizer-representable@cref}{{[originallemma][0][]00WS}{[1][0][]source}{}{}{}}
\makeatother
\title{Points of a site and points of its topos}
\author{The Stacks Project Authors}
\date{stacks-project:sites.tex:local-snapshot-20261012; GNU FDL 1.2 or later}
\begin{document}
\maketitle
\noindent\textit{Editorial note.} All delivered mathematical text is the author's own, copied from the pinned Stacks Project chapter file sites-cohomology.tex; no statement or proof step was rewritten, reordered or removed.\par\medskip
\noindent\textit{Editorial note.} The proof is the author's complete argument; the parenthetical 'verification omitted' at source line 8318 concerns the identification of two bijections on a representable sheaf, which the surrounding argument then uses and extends, and it is preserved verbatim inside the counted span.\par\medskip
\noindent\textit{Editorial note.} Every labelled reference occurring inside the counted spans is carried in assembly.external\_references as a non-load-bearing navigation pointer (same-chapter ones at their definition site, cross-chapter ones at the invocation inside the counted span), so no load-bearing mathematics is replaced by a stub. No bibliography citation occurs inside the counted spans.\par\medskip
\noindent\textit{Editorial note.} The checked chapter list for same-name collisions is the set of stacks chapters staged in this repository: chow, cohomology, divisors, etale-cohomology, flat, more-algebra, more-morphisms, morphisms, spaces-more-morphisms, varieties together with this card's chapter sites-cohomology.\par\medskip
\noindent\textit{Editorial note.} The preamble is the chapter's own retained minimal preamble (sources/stacks-sites-cohomology/retained-preamble.tex), a verbatim excerpt of the native preamble: the Xy-pic setup and the native theorem declarations with the mathematical macros.\par\medskip
\noindent\textit{Allowed standard prerequisites (editorial).}\begin{itemize}
\item The definition of a point of a site and of a topos. Same chapter, Definition definition-point and Definition definition-topos. Supply the conditions verified in both directions.
\item Coverings become surjections after sheafification. Same chapter, Lemma lemma-covering-surjective-after-sheafification. Supplies the surjectivity used to prove the first condition of a point of a site.
\item Sheafification commutes with finite limits. Same chapter, Lemma lemma-sheafification-exact. Supplies the commutation with fibre products used for the second condition.
\item Every sheaf is a coequalizer of coproducts of representables. Same chapter, Lemma lemma-sheaf-coequalizer-representable. Allows the comparison map to be checked on representable sheaves and then extended through colimits.
\end{itemize}\medskip
\ifdefined\hypertarget\hypertarget{sourcepointer-1}{}\fi
\noindent\textit{Source reference.} Original definition 00Y5: sources/stacks-sites/pinned-original-source.tex, source lines 8278--8278. Navigation pointer invoked at proof line 8278; the counted span contains the pointer only and the cited result is not reproduced inside it, so the pointer is registered as a non-load-bearing navigation stub in the established form for Stacks material. No bibliography citation is involved.\par\medskip
\ifdefined\hypertarget\hypertarget{sourcepointer-2}{}\fi
\noindent\textit{Source reference.} Original definition 00XA: sources/stacks-sites/pinned-original-source.tex, source lines 8280--8280. Navigation pointer invoked at proof line 8280; the counted span contains the pointer only and the cited result is not reproduced inside it, so the pointer is registered as a non-load-bearing navigation stub in the established form for Stacks material. No bibliography citation is involved.\par\medskip
\ifdefined\hypertarget\hypertarget{sourcepointer-3}{}\fi
\noindent\textit{Source reference.} Original lemma 00WT: sources/stacks-sites/pinned-original-source.tex, source lines 8286--8286. Navigation pointer invoked at proof line 8286; the counted span contains the pointer only and the cited result is not reproduced inside it, so the pointer is registered as a non-load-bearing navigation stub in the established form for Stacks material. No bibliography citation is involved.\par\medskip
\ifdefined\hypertarget\hypertarget{sourcepointer-4}{}\fi
\noindent\textit{Source reference.} Original lemma 00WJ: sources/stacks-sites/pinned-original-source.tex, source lines 8292--8292. Navigation pointer invoked at proof line 8292; the counted span contains the pointer only and the cited result is not reproduced inside it, so the pointer is registered as a non-load-bearing navigation stub in the established form for Stacks material. No bibliography citation is involved.\par\medskip
\ifdefined\hypertarget\hypertarget{sourcepointer-5}{}\fi
\noindent\textit{Source reference.} Original equation 04EH: sources/stacks-sites/pinned-original-source.tex, source lines 8301--8301. Navigation pointer invoked at proof line 8301; the counted span contains the pointer only and the cited result is not reproduced inside it, so the pointer is registered as a non-load-bearing navigation stub in the established form for Stacks material. No bibliography citation is involved.\par\medskip
\ifdefined\hypertarget\hypertarget{sourcepointer-6}{}\fi
\noindent\textit{Source reference.} Original lemma 00Y6: sources/stacks-sites/pinned-original-source.tex, source lines 8308--8308. Navigation pointer invoked at proof line 8308; the counted span contains the pointer only and the cited result is not reproduced inside it, so the pointer is registered as a non-load-bearing navigation stub in the established form for Stacks material. No bibliography citation is involved.\par\medskip
\ifdefined\hypertarget\hypertarget{sourcepointer-7}{}\fi
\noindent\textit{Source reference.} Original lemma 00Y8: sources/stacks-sites/pinned-original-source.tex, source lines 8235--8235. Navigation pointer invoked at proof line 8235; the counted span contains the pointer only and the cited result is not reproduced inside it, so the pointer is registered as a non-load-bearing navigation stub in the established form for Stacks material. No bibliography citation is involved.\par\medskip
\ifdefined\hypertarget\hypertarget{sourcepointer-8}{}\fi
\noindent\textit{Source reference.} Original lemma 00WS: sources/stacks-sites/pinned-original-source.tex, source lines 8321--8321. Navigation pointer invoked at proof line 8321; the counted span contains the pointer only and the cited result is not reproduced inside it, so the pointer is registered as a non-load-bearing navigation stub in the established form for Stacks material. No bibliography citation is involved.\par\medskip
\begin{lemma}
\label{lemma-point-site-topos}
Let $\mathcal{C}$ be a site.
\begin{enumerate}
\item Let $p$ be a point of the site $\mathcal{C}$.
Then the pair of functors $(p_*, p^{-1})$ introduced
above define a morphism of topoi
$\Sh(pt) \to \Sh(\mathcal{C})$.
\item Let $p = (p_*, p^{-1})$
be a point of the topos $\Sh(\mathcal{C})$.
Then the functor $u : U \mapsto p^{-1}(h_U^\#)$ gives
rise to a point $p'$ of the site $\mathcal{C}$
whose associated morphism of topoi $(p'_*, (p')^{-1})$
is equal to $p$.
\end{enumerate}
\end{lemma}

\begin{proof}
Proof of (1). By the above the functors $p_*$ and $p^{-1}$ are adjoint.
The functor $p^{-1}$ is required to be exact by
Definition \ref{definition-point}.
Hence the conditions imposed in
Definition \ref{definition-topos}
are all satisfied and we see that (1) holds.

\medskip\noindent
Proof of (2). Let $\{U_i \to U\}$ be a covering of $\mathcal{C}$.
Then $\coprod (h_{U_i})^\# \to h_U^\#$ is surjective, see
Lemma \ref{lemma-covering-surjective-after-sheafification}.
Since $p^{-1}$ is exact (by definition of a morphism of topoi) we conclude
that $\coprod u(U_i) \to u(U)$ is surjective.
This proves part (1) of
Definition \ref{definition-point}.
Sheafification is exact, see
Lemma \ref{lemma-sheafification-exact}.
Hence if $U \times_V W$ exists in $\mathcal{C}$, then
$$
h_{U \times_V W}^\# =  h_U^\# \times_{h_V^\#} h_W^\#
$$
and we see that $u(U \times_V W) = u(U) \times_{u(V)} u(W)$ since $p^{-1}$
is exact. This proves part (2) of
Definition \ref{definition-point}.
Let $p' = u$, and let $\mathcal{F}_{p'}$ be the stalk functor
defined by Equation (\ref{equation-stalk}) using $u$. There is
a canonical comparison map
$c : \mathcal{F}_{p'} \to \mathcal{F}_p = p^{-1}\mathcal{F}$.
Namely, given a triple $(U, x, \sigma)$ representing an element $\xi$
of $\mathcal{F}_{p'}$ we think of $\sigma$ as a map
$\sigma : h_U^\# \to \mathcal{F}$ and we can set
$c(\xi) = p^{-1}(\sigma)(x)$ since $x \in u(U) = p^{-1}(h_U^\#)$. By
Lemma \ref{lemma-points-recover}
we see that $(h_U)_{p'} = u(U)$. Since conditions (1) and (2) of
Definition \ref{definition-point}
hold for $p'$ we also have $(h_U^\#)_{p'} = (h_U)_{p'}$ by
Lemma \ref{lemma-point-pushforward-sheaf}.
Hence we have
$$
(h_U^\#)_{p'} = (h_U)_{p'} = u(U) = p^{-1}(h_U^\#)
$$
We claim this bijection equals the comparison map
$c : (h_U^\#)_{p'} \to p^{-1}(h_U^\#)$ (verification omitted).
Any sheaf on $\mathcal{C}$ is a coequalizer of maps
of coproducts of sheaves of the form $h_U^\#$, see
Lemma \ref{lemma-sheaf-coequalizer-representable}.
The stalk functor $\mathcal{F} \mapsto \mathcal{F}_{p'}$ and
the functor $p^{-1}$ commute with arbitrary colimits (as they
are both left adjoints).
We conclude $c$ is an isomorphism for every sheaf $\mathcal{F}$.
Thus the stalk functor $\mathcal{F} \mapsto \mathcal{F}_{p'}$
is isomorphic to $p^{-1}$ and we in particular see that it is exact.
This proves condition (3) of
Definition \ref{definition-point}
holds and $p'$ is a point. The final assertion has already been shown
above, since we saw that $p^{-1} = (p')^{-1}$.
\end{proof}

\end{document}
